\documentclass[11pt]{article}

\usepackage[margin=0.67in]{geometry}
\usepackage{setspace}
\usepackage{parskip}
\usepackage{hyperref}
\usepackage{titlesec}
\usepackage{xcolor}
\usepackage{array}
\usepackage{ragged2e}

\definecolor{PrimaryRed}{HTML}{8B1E1E}
\definecolor{SoftBlue}{HTML}{DDEEFF}
\definecolor{DarkGray}{HTML}{333333}

\hypersetup{
    colorlinks=true,
    linkcolor=black,
    urlcolor=black,
    citecolor=black
}

\titleformat{\section}{\large\bfseries\color{PrimaryRed}}{}{0em}{}
\titleformat{\subsection}{\normalsize\bfseries}{\thesubsection}{0.5em}{}

\setstretch{1.15}

\begin{document}

\begin{titlepage}
\newgeometry{margin=0.7in}
\thispagestyle{empty}

\begin{center}
    \vspace*{0.25in}
    {\large University of Rochester\\}
    {\normalsize Warner School of Education\\[0.35in]}

    {\color{PrimaryRed}\rule{\textwidth}{1.4pt}}\\[0.22in]
    {\Huge\bfseries Joy-Oriented Literacy Pursuit\\[0.12in]}
    {\Large Formal Reflective Journal\\[0.22in]}
    {\color{PrimaryRed}\rule{0.72\textwidth}{0.7pt}}\\[0.32in]

    \colorbox{SoftBlue}{%
        \parbox{0.86\textwidth}{%
            \centering\vspace{0.12in}
            \textbf{Classroom Technology, Design, and Literacy as Joy}\\[0.05in]
            A Teaching Neurodivergent Journal focused on Literature \& Joy via a Puzzle Plan
            \vspace{0.12in}
        }%
    }
\end{center}

\vspace{0.75in}

\begin{center}
\renewcommand{\arraystretch}{1.35}
\begin{tabular}{>{\bfseries\RaggedLeft}p{1.55in} p{4.85in}}
Student & Piter Garcia \\
Course & EDU498: Culturally-Historically Responsive Literacies \\
Class Session & Module 3 / Week 5 \\
Assignment & Joy-Oriented Literacy Pursuit \\
Context & Teaching Reflection through the Lens of this course \& a Puzzle Plan\\
Submission Date & June 05, 2026 \\
\end{tabular}
\end{center}

\vfill

\begin{center}
\begin{minipage}{0.7\textwidth}
\centering
\itshape
Joy-oriented literacy is not extra. It is part of access. When students feel capable, connected, and allowed to participate as themselves, literacy becomes a tool for liberation, self-determination, and empowerment.
\end{minipage}
\end{center}

\restoregeometry
\end{titlepage}
\clearpage

\doublespacing

\section*{What I Created and Who I Shared It With}

For this joy-oriented literacy pursuit, I designed and facilitated technology lessons where students experienced literacy through building, collaboration, speaking, movement, and problem-solving. The most joyful examples came from my Kindergarten circuit and airplane-building lessons and my Minecraft lessons in third and fourth grade. These activities were shared with the classroom community because students built, watched, helped, responded, revised, and learned from one another in real time.

In my Kindergarten circuit lessons, students gathered around a shared demonstration, named parts, held materials, and connected pieces before counting down together to turn on a motor-based fan. They then worked in teams to build an airplane, assigning roles like wings, body, wheels, and fan connection, so that every student had a meaningful way to participate. In my fourth grade Minecraft lessons, students who finished their work were invited to help classmates who were not done. That shift turned task completion into peer teaching and community responsibility, and made the shared experience itself a literacy product: students became makers, helpers, designers, and explainers.

\section*{My Creative and Learning Process}

My process began with a question: \textit{how can a technology lesson become more than task completion?} Knobel and Lankshear (2014) helped me answer it: literacy is not fixed to a single form; it evolves as tools, communities, and purposes change. Price-Dennis, Holmes, and Smith (2015) pushed further by showing that digital tools only become inclusive when teachers intentionally use them to reduce barriers, increase participation, and support meaningful communication. Together, their work reframed my lessons as genuine literacy work because students were interpreting systems, building vocabulary, communicating with peers, and representing thinking through more than one mode.

In the circuit and airplane activities, students connected vocabulary to action. They did not only hear the word ``switch''; they held it, connected it, and watched what happened when it worked. In the Grade 5 airplane and fan-circuit challenge, students were assigned roles like building wings, shaping the body, making wheels, and connecting the fan to the switch, which made the work collaborative and distributed rather than individual. In Minecraft, inviting early finishers to help peers shifted the room from ``I am done'' to ``I can help someone else succeed,'' where one student troubleshoots while another consolidates mastery.

Joy was not an accident of fun activities. It happened because students had access to materials, roles, language, peer support, and visible success. This matters especially for students who are multilingual, neurodivergent, or whose ways of knowing are not typically recognized by traditional literacy tasks. Knobel and Lankshear describe new literacies as social practices built around participation and community, not just technical use of tools. Price-Dennis and colleagues argue that adding a device to a traditional task is not innovation; the real question is whether the task creates more access, voice, and critical thinking. That became the design principle: \textit{enough structure to feel safe, enough freedom to feel ownership}.

\section*{How This Cultivated Joy, Love, and Aesthetic Fulfillment}

Joy came from students seeing their work become real. In the Kindergarten fan-circuit lessons, the class counted down together and watched the fan turn on. The science was no longer abstract; it was visible, audible, and shared. Aesthetic fulfillment came from building something with shape, motion, and purpose: their work had form, effort, and a relationship to the physical world. The joy was not only in the fan working; it was in the class experiencing the success together.

Love came through relationships. In the fourth grade Minecraft lessons, students who understood the task helped classmates who were still working. That kind of peer support changed the emotional meaning of the activity. Success was not private or competitive; it became something students could give each other. For students who often experience school as a place where their pace, language, body, or attention is a problem, that moment of being the one who helps carries a different kind of meaning.

Joy in literacy does not always look quiet or neat. Sometimes it looks like a student trying again, a group arguing productively over what goes where, a fan finally turning on, a peer stepping in to help, or a student realizing that their hands, ideas, and voice matter. That is aesthetic fulfillment: learning feels alive, students can see themselves inside the work, and the classroom becomes a place where knowledge is built rather than delivered.

\section*{Muhammad's Framework and the Case for Literacy as Access}

Muhammad's framework asks that a literacy pursuit develop identity, skills, intellect, criticality, and joy together. This project did not aim only at science skills; it asked students to see themselves as builders, helpers, problem-solvers, and contributors, because literacy should help students understand themselves, participate in community, and experience agency. Through the Puzzle Plan lens, that matters especially for students whose identities, languages, or learning needs are often treated as deficits rather than resources.

Knobel and Lankshear (2014) gave me language for what I was already seeing: building, explaining, and peer teaching are new literacies because meaning-making happens through digital, visual, participatory, and networked forms, not only through print. Their work challenges the assumption that school-approved print performance is the only valid measure of intelligence or literacy. Price-Dennis, Holmes, and Smith (2015) connect this to inclusion: when tasks allow for multimodal creation, oral explanation, and collaborative design, more students can show what they know, including students who are neurodivergent, multilingual, or differently abled.

Brooks and Frankel (2020) matter here because students do not arrive as blank technology users. Their work on transnational digital literacy shows that students may already be navigating languages, countries, and communities through digital tools before school names any of it academic. When a student builds in Minecraft, explains a circuit orally, helps a peer troubleshoot, or draws a design, that student is using literacy in a way that deserves recognition. These are not side activities; they require interpretation, audience awareness, cultural knowledge, and digital navigation.

Melo-Pfeifer and Gertz (2022) pushed me toward criticality: digital spaces are not neutral, and students need tools to question headlines, missing voices, representation, and assumptions, especially students whose communities are most often misrepresented. The Common Sense Digital Citizenship materials extend this by giving students language for privacy, confirmation bias, online identity, and upstander communication. My Puzzle Plan builds exactly those habits: collaboration, questioning, audience awareness, and responsibility to others.

Stornaiuolo and Thomas (2017) reframe students as active participants who use digital tools to speak and create public meaning, not passive consumers. Pinkard's ecosystem view reinforces this: students need tools, mentors, audiences, and real reasons to create. The Grade 4 peer-teaching moment was that ecosystem in action. Serpagli and Mensah (2021) connect to my placement directly: students explained technical ideas through familiar, social, and active forms. Across all these readings, the argument is the same: digital literacy should position students as creators with voice, identity, and purpose.

\section*{How This Informs My Future Teaching}

Literacy is not limited to written language. It is also building, coding, designing, explaining, translating, troubleshooting, and helping. In my classroom, I will design activities where students have multiple entry points and different ways to show their understanding. This is especially important for students who may be multilingual, neurodivergent, disabled, chronically ill, culturally diverse, or simply different from the ``default'' learner that schools often imagine.

I will keep designing for authentic audiences: Minecraft worlds tied to a science concept, bilingual or multimodal explanations, collaborative classroom murals, peer-created tutorials, or circuitry projects explained through speaking, drawing, writing, or video. Peer teaching will be built into the culture because students often understand each other's barriers in ways adults miss, and because being the one who helps is itself a form of literacy and recognition.

The core question my Puzzle Plan and EQUITAS lens keeps asking is not whether a task looks traditional enough to count as literacy. It is whether the task gives students real access to meaning, identity, criticality, and agency. Knobel and Lankshear show that literacy changes with communities and tools; Price-Dennis and colleagues show that inclusion is designed, not assumed; Brooks and Frankel show that students' existing practices already are literacy; and Stornaiuolo and Thomas show that digital tools can be instruments of voice, not just task completion. When students feel capable, connected, and allowed to participate as themselves, literacy becomes a tool for liberation, self-determination, and empowerment.

\end{document}
